\label{sec:conclusions}
For $\alpha=1/2$ at $D=1$~Gpc with $A_\alpha=1$:
$\sigma_\ell = \sqrt{\ell_{\rm P}\,D} \approx 2.2\times10^{-5}$~m,
giving $\sigma_\phi \approx 280$ (fringes completely erased at $\sigma_\phi\gg 1$),
while $c\tau_c = \lambda^2/(2\pi\Delta\lambda) \approx 80\,\mu$m and
$\sigma_\Delta/\tau_c \approx \sqrt{2}\,\sigma_\ell/(c\tau_c) \approx 0.40$.
The narrow-band expansion parameter is
$\delta t_{\rm rms}/\tau_c = \sigma_\ell/(c\tau_c) \approx 0.28$,
so corrections to Eq.~\eqref{eq:g2unchanged} are at the $\sim 8\%$ level;
the exact convolution~\eqref{eq:g2obs} should be used for precision work.
Critically, this case ($\alpha=1/2$) is already excluded by fringe-visibility
measurements~\cite{Perlman2015,Christiansen2011,2003ApJ...587L...1R}.
For the holographic model $\alpha=2/3$ with the same parameters, $\sigma_\phi\sim 10^{-8}$,
so all foam coherence observables are negligibly small with current instrumentation.
The theoretical framework is therefore most useful as a basis for future constraints
at shorter wavelengths (X-ray or near-UV), longer baselines, or next-generation
intensity interferometers, where the sensitivity scales as $\lambda^{-1}$ and $D^{1-\alpha}$.
To make this concrete: for $\alpha=2/3$ at $D=10$~Gpc and soft X-ray wavelength
$\lambda=1$~\AA, one finds $\sigma_\ell \approx 4.3\times10^{-15}$~m,
$\sigma_\phi \approx 2.7\times10^{-4}$, and
$\sigma_\Delta/\tau_c \approx 3.8\times10^{-7}$ (for $\Delta\lambda/\lambda=10^{-3}$) —
still approximately seven orders of magnitude below a unity HBT signal.
The framework therefore serves primarily as a theoretical tool and as an upper-limit
setter for future precision measurements, rather than as a near-term positive-detection
strategy.
